\documentclass{article}

\usepackage{booktabs}
\usepackage{tabularx}

\title{Development Plan\\\progname}

\author{\authname}

\date{}

\input{../Comments.text}
\input{../Common.text}

\begin{document}

\maketitle

\begin{table}[hp]
\caption{Revision History} \label{TblRevisionHistory}
\begin{tabularx}{\textwidth}{llX}
\toprule
\textbf{Date} & \textbf{Developer(s)} & \textbf{Change}\\
\midrule
26 September 2026 & Moustafa Moustafa & Added my reflection responses to the appendix.\\
26 September 2026 & Rayan Nasrallah & Added reflection responses on the value
of a development plan, the advantages and disadvantages of CI/CD, and the
team's disagreement over naming conventions.\\
\bottomrule
\end{tabularx}
\end{table}

\newpage{}

\wss{Put your introductory blurb here.  Often the blurb is a brief roadmap of
what is contained in the report.}

\wss{Additional information on the development plan can be found in the
\href{https://gitlab.cas.mcmaster.ca/courses/capstone/-/blob/main/Lectures/L02b_POCAndDevPlan/POCAndDevPlan.pdf?ref_type=heads}
{lecture slides}.}

\section{Confidential Information?}

\wss{State whether your project has confidential information from industry, or
not.  If there is confidential information, point to the agreement you have in
place.}

\wss{For most teams this section will just state that there is no confidential
information to protect.}
\section{IP to Protect}

\wss{State whether there is IP to protect.  If there is, point to the agreement.
All students who are working on a project that requires an IP agreement are also
required to sign the ``Intellectual Property Guide Acknowledgement.''}

\section{Copyright License}

\wss{What copyright license is your team adopting.  Point to the license in your
repo.}

\section{Team Meeting Plan}

\wss{This section is for meetings between the team members, and meetings between 
the team and their supervisor/advisor (as appropriate).}

\wss{How often will the team meet? where? when?}

\wss{If the meeting is a physical location (not virtual), out of an abundance of
caution for safety reasons you shouldn't put the location online}

\wss{To help schedule meetings, you can use the lecture/tutorial time when we do 
not have anything else scheduled.}

\wss{How often will you meet with your supervisor/advisor?  when?  where?}

\wss{Will meetings be virtual?  At least some meetings should likely be
in-person.}

\wss{How will the meetings be structured?  There should be a chair for all meetings.  
There should be an agenda for all meetings.}

\section{Team Communication Plan}

\wss{Issues on GitHub should be part of your communication plan.}

\section{Team Member Roles}

\wss{You should identify the types of roles you anticipate, like notetaker,
leader, meeting chair, reviewer.  Assigning specific people to those roles is
not necessary at this stage.  In a student team the role of the individuals will
likely change throughout the year.}

\section{Workflow Plan}

\begin{itemize}
	\item How will you be using git, including branches, pull request, etc.?
	\item How will you be managing issues, including template issues, issue
	classification, etc.?
  \item Use of CI/CD
\end{itemize}

\section{Project Decomposition and Scheduling}

\begin{itemize}
  \item How will you be using GitHub projects?
  \item Include a link to your GitHub project
\end{itemize}

\wss{How will the project be scheduled?  This is the big picture schedule, not
details. You will need to reproduce information that is in the course outline
for deadlines.}

\section{Proof of Concept Demonstration Plan}

What is the main risk, or risks, for the success of your project?  What will you
demonstrate during your proof of concept demonstration to convince yourself that
you will be able to overcome this risk?

\section{Expected Technology}

\wss{What programming language or languages do you expect to use?  What external
libraries?  What frameworks?  What technologies.  Are there major components of
the implementation that you expect you will implement, despite the existence of
libraries that provide the required functionality.  For projects with machine
learning, will you use pre-trained models, or be training your own model?  }

\wss{The implementation decisions can, and likely will, change over the course
of the project.  The initial documentation should be written in an abstract way;
it should be agnostic of the implementation choices, unless the implementation
choices are project constraints.  However, recording our initial thoughts on
implementation helps understand the challenge level and feasibility of a
project.  It may also help with early identification of areas where project
members will need to augment their training.}

Topics to discuss include the following:

\begin{itemize}
\item Specific programming language
\item Specific libraries
\item Pre-trained models
\item Specific linter tool (if appropriate)
\item Specific unit testing framework
\item Investigation of code coverage measuring tools
\item Specific plans for Continuous Integration (CI), or an explanation that CI
  is not being done
\item Specific performance measuring tools (like Valgrind), if
  appropriate
\item Tools you will likely be using?
\end{itemize}

\wss{git, GitHub and GitHub projects should be part of your technology.}

\section{Use of AI and LLMs}

\wss{This section is a continuation of the previous section.  The role of LLMs
in your project as important enough to warrant their own section.  
How will your team be using LLMs for code and documentation development?  
What tools will you be using?  How will you verify the output of your LLMs?}

\section{Coding Standard}

\wss{What coding standard will you adopt?}

\newpage{}

\section{Appendix --- Generative AI Use Disclosure}

\wss{ChatGPT (version 5) was used to brainstorm ideas for the template for this disclose section.}

\wss{This section is about the use of AI to write this document, not your
planned use of AI for your project.  You discuss that topic in one of the
sections in the body of this report.}

\subsection{AI Tools Used}

\wss{Give the name of the AI tool(s) used and the version.}

\subsection{How and Where Used}

\wss{Explain what the tool(s) were used for.  Examples include: brainstorming,
explaining a technical concept, reviewing the document, generating or modifying
text, generating or modifying diagrams, identifying errors or omissions or
inconsistencies.}

\wss{For each of the uses, identify which parts of the document were affected.}

\wss{You do not need to report every interaction, but you should disclose the 
uses that materially contributed to your work}

\subsection{Verification of AI-Generated Content}

\wss{\begin{itemize}
    \item Describe how the team checked AI-generated or AI-assisted material
    before including it in the document.
    \item In particular, verify factual claims, technical assertions,
    requirements, calculations, references, and other information for which an
    AI system may produce incorrect or fabricated results.
    \item \textbf{The team is responsible for the accuracy, completeness,
consistency, and appropriateness of the submitted document, regardless of
whether material was generated by a human or an AI system.}
\end{itemize}}

\newpage{}

\section*{Appendix --- Reflection}

\input{../Reflection.text}


\begin{enumerate}

    \item \textbf{Why is it important to create a development plan
    prior to starting the project?}

    \textbf{Ahmed Elzaria:}

    Creating the development plan helped me recognize that agreeing
    on a project idea does not automatically mean we agree on how
    to complete it. A team can share the same goals while having
    different expectations about responsibilities, deadlines, and
    acceptable contributions. Discussing these expectations early
    gives us a common basis for coordinating work and addressing
    problems before they affect a major deliverable.

    While developing the team roles, I initially considered
    assigning each person a narrow area of implementation.
    However, I realized that this could restrict members'
    learning and create dependence on individuals. Making every
    responsibility shared introduced the opposite problem:
    accountability became unclear. The approach I settled on
    combines named leads with shared development duties and
    explicit task ownership. This showed me that planning
    involves balancing competing needs, rather than simply
    dividing a list of tasks.

    The technology discussion also demonstrated the value of
    planning. Clarifying that enhanced drawings should remain
    editable changed what we needed from the implementation:
    returning a generated image would not satisfy that
    expectation. Identifying this distinction early helped
    expose a feasibility question before investing in an
    unsuitable approach. I therefore see the development plan
    as a starting point that makes assumptions visible and
    guides early investigation, while remaining open to revision
    as we learn.

    \textbf{Moustafa Moustafa:}

    A development plan makes it clear who is responsible for what
    before the work starts. Without one, tasks are unclear and it
    is easy for someone to say ``I don't want to do this'' with
    nothing to point back to. When we define roles, who works on
    what, how we split tasks, and which tools each person uses,
    it saves a lot of friction later.

    I have learned this from experience. On a previous project, we
    did not have a well-defined plan, and it came back to haunt us
    near the end. All three of us got busy with university just as
    some milestones were due. Nobody wanted to take them on, but
    someone had to, and the work ended up being split unevenly accross 2 people.
    With a clear plan, clearly separated tasks, and agreed
    roles, we could have avoided that situation entirely. That is
    why I think having the plan in place from the start matters so
    much for a five-person capstone team.

    \textbf{Rayan Nasrallah:}

    The benefit I found clearest was being made to sketch the
    technical implementation at a high level, because doing so
    exposed what we did not yet know how to build. Committing to a
    rough shape for the system forced questions that a feature list
    never raises on its own.
    Our discussion about keeping refined drawings editable is a good
    example: the requirement sounds straightforward until you
    consider the implementation, at which point returning a
    generated image clearly fails to satisfy it. Identifying that
    early gives us time to investigate, rather than discovering the
    problem once work depends on the wrong approach.

    These feasibility questions feed directly into scope, which
    matters more here than it would on an open-ended project. We
    have two semesters, five members, and a fixed set of deliverables
    alongside the implementation, so the amount of work a feature
    requires is as important as how valuable it is. Thinking this
    through in advance is what let us treat structured practice,
    adaptive guidance, and progress reporting as core work while
    moving camera-based tracing and letter activities to stretch
    goals. Without estimating the effort first, that judgement would
    have been guesswork.

    I also think there is a real difference between discussing these
    decisions and writing them down. Discussion can produce the
    feeling of agreement while members leave with different
    understandings of what was decided, and writing forces the
    vagueness to surface. The written plan also gives us a concrete
    reference to revise. I do expect parts of it to prove wrong once
    implementation begins, since we are planning around technology
    we have not used yet, so I see its value less in being accurate
    and more in making our assumptions explicit enough to correct.

    \item \textbf{In your opinion, what are the advantages and
    disadvantages of using CI/CD?}

    \textbf{Ahmed Elzaria:}

    I see the main advantage of Continuous Integration (CI) as
    giving the team consistent feedback whenever changes are
    proposed. With five members contributing, relying on each
    person to remember and reproduce every check would make
    mistakes easier to miss. Automated builds, linting, and tests
    can identify problems while the relevant change is still
    small and easier to investigate. For our documentation, an
    automated \LaTeX{} build can also reveal compilation errors
    before submission.

    Continuous Delivery or Deployment (CD) extends this process.
    Continuous delivery keeps changes ready for release, while
    continuous deployment automatically releases changes that
    pass the required checks. For Squigglr, preview deployments
    would make it easier for teammates to review an interface
    change and try it on a tablet. However, passing automated
    checks should not be treated as proof that a change is ready
    for users. Drawing responsiveness and stylus behaviour still
    need evaluation on actual devices, and a successful document
    build does not establish that its content satisfies the
    rubric.

    The main disadvantages are the effort required to configure
    and maintain the workflows, the time spent investigating
    failed runs, and the risk of trusting incomplete tests too
    much. A pipeline can pass while important behaviour remains
    untested. For example, a scoring function could produce
    consistently incorrect results if its tests use the same
    mistaken assumptions as its implementation.

    Given our timeline, I favour introducing a small set of useful
    checks first and expanding them as the application develops.
    This would allow us to gain early feedback without spending
    disproportionate effort on automation before the core
    functionality exists. Human review and device testing would
    remain part of our process.

    \textbf{Moustafa Moustafa:}

    I first saw the real value of CI/CD during my previous co-op. It was
    my first experience on a large project, and I realized that a
    codebase of that size could not function without it once more
    than a few people were contributing. The same applies on a
    smaller scale to capstone. With five people working in the same
    repository, we would constantly run into conflicts and broken
    builds without CI/CD, and it would be hard to tell which change
    caused a problem. Automated checks on every pull request catch
    issues early, make it clear which change introduced them, and
    make it much easier to roll back when something goes wrong.

    The main disadvantage is that it can become a bit of a rabbit hole. Some
    people keep improving their pipeline, adding more optimizations
    and tools, until maintaining the CI/CD takes more time than the
    actual work. There is a balance to find: the pipeline should
    be good enough to let the team work together safely and recover
    from mistakes, without being over-engineered. For us, that means
    keeping CI/CD simple and only adding to it when we actually need
    to.

    \textbf{Rayan Nasrallah:}

    Continuous Integration means merging our work into the shared
    branch frequently, with automated checks running on it each
    time, instead of integrating large amounts of work at once.
    Continuous Delivery keeps every change that passes those checks
    ready to release, while continuous deployment releases it
    without waiting for a person. Our repository already does all
    of this for documentation: a GitHub Actions workflow determines
    which documents a change affects, compiles them, commits the
    resulting PDFs, and publishes them to our project page. Once
    implementation begins, the same pipeline would run tests and
    linting rather than only document builds.

    The advantage I value most is the effect this has on how we
    work, rather than what it automates. Because every change
    arrives as a pull request that has to build before it can
    merge, work cannot sit unintegrated on a branch for weeks and
    then arrive all at once near a deadline. Changes stay small
    enough for somebody other than the author to read properly, and
    the project advances in increments we can each follow.

    The published output stays current for the same reason. Our
    PDFs always correspond to what is on the main branch, so nobody
    has to compile everything before a submission and discover a
    problem then. For the application, preview deployments would
    let the rest of the team try an interface change on a tablet
    instead of taking the author's description of it on trust.

    I initially thought there were no real disadvantages, and
    looking more carefully at our own workflow changed that view.
    The clearest cost is that a passing pipeline is easy to mistake
    for a correct result. Our problem statement compiled without
    error while citing references the document did not yet contain,
    because \LaTeX{} has no reason to object to that. The same
    failure will be available to us in code, where tests can pass
    while drawing responsiveness on a real tablet remains wrong.
    Automated checks only cover what we thought to check, so I read
    a passing run as evidence that nothing obvious broke rather
    than that the work is right.

    \item \textbf{What disagreements did your group have in this
    deliverable, if any, and how did you resolve them?}

    \textbf{Ahmed Elzaria:}

    Our main differences concerned the consequences of missing meetings
    and the roles and expectations of each team member. These discussions
    required more communication than simply dividing the deliverable
    because they established how we would hold one another accountable
    throughout the project.

    For attendance and missed commitments, I wanted expectations that
    were clear enough to act on while recognizing that illness,
    emergencies, and unavoidable conflicts can occur. I saw a risk in
    both extremes: vague consequences could allow repeated problems
    to continue, while an overly rigid response could treat a genuine
    difficulty as a lack of effort. Through discussion, we developed
    an approach with measurable triggers and a staged response.
    The charter begins with a private conversation and a recovery
    plan, followed by escalation to the teaching team if commitments
    remain unmet without an acceptable explanation. This gives the
    team a basis for taking action while allowing us to understand
    the circumstances first.

    We also differed over how specifically to define individual roles.
    While working on this section, I found that narrow assignments
    could make members feel restricted to one component, whereas
    making everything a shared responsibility left no clear person
    to coordinate the work. The resulting structure assigns broad
    lead roles while making software development, testing,
    documentation, and peer review responsibilities of every member.
    Assigning Moustafa as the notetaker provides a consistent owner
    for meeting records, while rotating the quality assurance
    coordinator each milestone distributes deliverable review duties.

    These discussions changed how I understood fairness within the
    team. Giving everyone the same number of titles does not
    necessarily create an equal workload, because implementation,
    coordination, and review require different amounts of effort
    at different stages. I now see the role assignments as a
    starting point that must be supported by clear task ownership
    and regular workload reviews.

    My main lesson was that expectations need to be discussed
    explicitly rather than assumed. Writing down the consequences,
    role boundaries, and shared duties helped turn our differing
    views into a practical agreement. However, agreement on the
    document does not guarantee that the process will work well.
    We will need to revisit it as the workload increases and check
    whether the responsibilities and responses remain fair and useful.

    \textbf{Rayan Nasrallah:}

    Our clearest disagreement was about naming conventions for
    branches and commit messages. It is a small matter next to the
    project itself, which is part of why it took as long as it did:
    nobody treated it as urgent, but everybody touches it on every
    piece of work, so the inconsistency kept reappearing.

    The first split was over whether a branch and commit should
    reference the issue they belong to. Moustafa and I both thought
    they should, since it is what lets somebody move from a commit
    back to the discussion that produced it without asking the
    author. The rest of the team did not see it as significant
    either way, and I do not think they were being unreasonable,
    because on a five-person project you can usually work out who
    did what by asking. What convinced them was how much easier it
    makes the history to move around in. With the issue number in
    the branch and the commit, searching GitHub for it brings up the
    pull request, the review comments, and the discussion together,
    so anybody can follow a change back to its reasoning on their
    own. Nothing requires us to work this way, but the alternative
    is that the connection lives only in whoever happens to
    remember it.

    Agreeing on that did not settle the details. I wanted
    underscores in branch names and the full branch name repeated in
    the commit message, and the team disagreed with both. Moustafa
    argued for dashes instead of underscores, which we accepted, and
    Ahmed used that to propose the pattern we now follow, combining
    a type, the issue number, and a dash-separated description.
    Ahmed and Moustafa then each proposed a different format for
    commit messages, and we settled on the one that carries the
    issue number in the subject line rather than in a trailer at the
    end.

    What resolved it in the end was writing it down. The
    back-and-forth continued while the convention existed only in
    conversation, because each of us remembered the agreement
    slightly differently, and it stopped once we recorded the
    branch and commit formats in the repository README where anybody
    can check them.

\end{enumerate}

\newpage{}

\section*{Appendix --- Team Charter}

\wss{borrows from
\href{https://engineering.up.edu/industry_partnerships/files/team-charter.pdf}
{University of Portland Team Charter}}

\subsection*{External Goals}

\wss{What are your team's external goals for this project? These are not the
goals related to the functionality or quality fo the project.  These are the
goals on what the team wishes to achieve with the project.  Potential goals are
to win a prize at the Capstone EXPO, or to have something to talk about in
interviews, or to get an A+, etc.}

\subsection*{Attendance}

\subsubsection*{Expectations}

\wss{What are your team's expectations regarding meeting attendance (being on
time, leaving early, missing meetings, etc.)?}

\subsubsection*{Acceptable Excuse}

\wss{What constitutes an acceptable excuse for missing a meeting or a deadline?
What types of excuses will not be considered acceptable?}

\subsubsection*{In Case of Emergency}

\wss{What process will team members follow if they have an emergency and cannot
attend a team meeting or complete their individual work promised for a team
deliverable?}

\subsection*{Accountability and Teamwork}

\subsubsection*{Quality} 

\wss{What are your team's expectations regarding the quality
of team members' preparation for team meetings and the quality of the
deliverables that members bring to the team?}

\subsubsection*{Attitude}

\wss{What are your team's expectations regarding team members' ideas,
interactions with the team, cooperation, attitudes, and anything else regarding
team member contributions?  Do you want to introduce a code of conduct?  Do you
want a conflict resolution plan?  Can adopt existing codes of conduct.}

\subsubsection*{Stay on Track}

\wss{What methods will be used to keep the team on track? How will your team
ensure that members contribute as expected to the team and that the team
performs as expected? How will your team reward members who do well and manage
members whose performance is below expectations?  What are the consequences for
someone not contributing their fair share?}

\wss{You may wish to use the project management metrics collected for the TA and
instructor for this.}

\wss{You can set target metrics for attendance, commits, etc.  What are the
consequences if someone doesn't hit their targets?  Do they need to bring the
coffee to the next team meeting?  Does the team need to make an appointment with
their TA, or the instructor?  Are there incentives for reaching targets early?}

\subsubsection*{Team Building}

\wss{How will you build team cohesion (fun time, group rituals, etc.)? }

\subsubsection*{Decision Making} 

\wss{How will you make decisions in your group? Consensus?  Vote? How will you
handle disagreements? }

\end{document}